\chapter{Risque, validation et réglementation}\label{ch:iv2:risk-validation-and-regulation}

\begin{ivtrains}
La valeur en risque (VaR) et l'expected shortfall, calculées à la main et lorsqu'elles se brisent, les backtests et ce qu'ils permettent réellement de savoir, les tests de résistance, le risque factoriel, les règles de capital auxquelles une réponse sur le risque doit résister, et la validation d'un modèle. Tous les portefeuilles présentés ici sont inventés. Quantiles de la loi normale~: 2.326 à 99~\% et 1.96 à 97.5~\%.
\end{ivtrains}

\begin{question}[$\star$ \iqroles{Risque, QR, Trader} \iqfirm{toutes} \iqlevel{junior}]\label{iq:iv2:risk-validation-and-regulation:1}
Vous êtes long de 100 millions de dollars d'un indice actions dont la volatilité quotidienne est de 1~\%. Quelle est la valeur en risque (VaR) à un jour à 99~\% sous un modèle normal de moyenne nulle~? Que donne la règle de la racine carrée du temps pour dix jours~?
\end{question}

\begin{question}[$\star$ \iqroles{Risque, QR, Analyste} \iqfirm{toutes} \iqlevel{junior}]\label{iq:iv2:risk-validation-and-regulation:2}
Vous détenez 50 millions de dollars d'un indice dont la volatilité quotidienne est de 1~\%, et 30 millions de dollars d'un autre dont la volatilité quotidienne est de 1.5~\%. Leur corrélation est de 0.8. Quelle est la VaR à un jour à 99~\%~? Combien la diversification permet-elle d'économiser~?
\end{question}

\begin{question}[$\star$ \iqroles{Risque, QR} \iqfirm{toutes} \iqlevel{junior}]\label{iq:iv2:risk-validation-and-regulation:3}
Définissez l'expected shortfall. Pour une loi normale de moyenne 0 et d'écart type 1, quelle est sa valeur à 97.5~\% et à 99~\%~? Comment la valeur à 97.5~\% se compare-t-elle à la VaR à 99~\%~?
\end{question}

\begin{question}[$\star$ \iqroles{Risque, Analyste, QR} \iqfirm{toutes} \iqlevel{junior}]\label{iq:iv2:risk-validation-and-regulation:4}
Sur 250 jours, vos cinq pires pertes quotidiennes, en millions, ont été de 8, 6, 5, 4 et 3.5. Estimez la VaR historique à 99~\% et l'expected shortfall à 99~\%. Indiquez la convention que vous retenez.
\end{question}

\begin{question}[$\star$ \iqroles{Risque, QR} \iqfirm{banque} \iqlevel{junior}]\label{iq:iv2:risk-validation-and-regulation:5}
La VaR à un jour à 99~\% d'une banque a enregistré 7 dépassements au cours des 250 derniers jours. Selon le feu tricolore de Bâle, dans quelle zone cela tombe-t-il, et quel effet cela a-t-il sur le multiplicateur de 3~?
\end{question}

\begin{question}[$\star$ \iqroles{Risque, QR, QT} \iqfirm{toutes} \iqlevel{junior}]\label{iq:iv2:risk-validation-and-regulation:6}
Une volatilité à pondération exponentielle utilise $\lambda = 0.94$. Quelle est la demi-vie du poids d'un rendement, en jours~? Hier, l'estimation de la volatilité quotidienne était de 1~\%, et le rendement d'aujourd'hui est de 2~\%. Quelle est la nouvelle estimation~?
\end{question}

\begin{question}[$\star$ \iqroles{Risque, Analyste, QR} \iqfirm{toutes} \iqlevel{junior}]\label{iq:iv2:risk-validation-and-regulation:7}
Une action a un bêta de 1.2 par rapport à un marché dont la volatilité est de 15~\%, et une volatilité idiosyncratique de 20~\%. Quelles sont sa volatilité totale et son $R^2$~? Quelle est la volatilité d'un portefeuille à pondérations égales de 100 actions de ce type, dont les parties idiosyncratiques sont indépendantes~?
\end{question}

\begin{question}[$\star$ \iqroles{Risque, Dev, QD} \iqfirm{toutes} \iqlevel{junior}]\label{iq:iv2:risk-validation-and-regulation:8}
La valeur d'un compte évolue ainsi~: 100, 120, 90, 130, 65, 80. Quel est le drawdown maximal~? Comment le calculez-vous en un seul parcours~?
\end{question}

\begin{question}[$\star$ \iqroles{Risque, Trader, QT} \iqfirm{toutes} \iqlevel{junior}]\label{iq:iv2:risk-validation-and-regulation:9}
Vous êtes long d'une option d'achat de delta 0.5 et de gamma 0.02 sur une action cotée 100 dollars, dont la volatilité quotidienne est de 2~\%. Quelle est la VaR à un jour à 99~\% delta-normale par option~? Est-elle trop élevée ou trop faible une fois le gamma pris en compte~?
\end{question}

\begin{question}[$\star$ \iqroles{Risque, Trader, Analyste} \iqfirm{banque, hedge fund} \iqlevel{junior}]\label{iq:iv2:risk-validation-and-regulation:10}
Un portefeuille obligataire présente un DV01 de 50\,000 dollars, et le rendement a une volatilité quotidienne de 6 points de base. Quelle est la VaR à un jour à 99~\%~?
\end{question}

\begin{question}[$\star\star$ \iqroles{Risque, QR} \iqfirm{toutes} \iqlevel{junior, confirmé}]\label{iq:iv2:risk-validation-and-regulation:11}
Deux obligations indépendantes perdent chacune 100 avec une probabilité de 4~\%, et rien sinon. Quelles sont la VaR à 95~\% de chacune d'elles et celle de la paire~? Comparez l'expected shortfall à 95~\%. Quelle propriété cela met-il en évidence que la VaR ne possède pas~?
\end{question}

\begin{question}[$\star\star$ \iqroles{Risque, QR} \iqfirm{banque} \iqlevel{junior, confirmé}]\label{iq:iv2:risk-validation-and-regulation:12}
Votre VaR à 99~\% a connu 7 dépassements en 250 jours. Effectuez le test de Kupiec. Rejetez-vous à 5~\%~? Quelles sont la probabilité de queue binomiale exacte et la conclusion du test pour 4 dépassements, puis pour 0~?
\end{question}

\begin{question}[$\star\star$ \iqroles{Risque, QR} \iqfirm{banque} \iqlevel{junior, confirmé}]\label{iq:iv2:risk-validation-and-regulation:13}
Sur 250 paires de jours, une VaR à 99~\% présente les effectifs suivants~: aucun dépassement puis aucun dépassement, 238~; aucun dépassement puis dépassement, 5~; dépassement puis aucun dépassement, 4~; dépassement puis dépassement, 3. Calculez les probabilités de transition et effectuez le test d'indépendance de Christoffersen.
\end{question}

\begin{question}[$\star\star$ \iqroles{Risque, QR} \iqfirm{toutes} \iqlevel{junior, confirmé}]\label{iq:iv2:risk-validation-and-regulation:14}
Les rendements quotidiens présentent une autocorrélation d'ordre 1 de 0.2, selon un AR(1). De combien la règle de la racine carrée du temps sous-estime-t-elle la volatilité à dix jours~? Qu'advient-il avec une autocorrélation négative~?
\end{question}

\begin{question}[$\star\star$ \iqroles{Risque, QR} \iqfirm{toutes} \iqlevel{junior, confirmé}]\label{iq:iv2:risk-validation-and-regulation:15}
Les rendements suivent une loi de Student-$t$ à 4 degrés de liberté, réajustée à une variance de 1. Le quantile à 1~\% de la loi $t_4$ non réajustée vaut $-3.747$. Quelles sont la VaR et l'expected shortfall à 99~\%, comparées à celles de la loi normale~?
\end{question}

\begin{question}[$\star\star$ \iqroles{Risque, QR, QD} \iqfirm{toutes} \iqlevel{junior, confirmé}]\label{iq:iv2:risk-validation-and-regulation:16}
Vous estimez une VaR à 99~\% par Monte Carlo sur un P\&L normal à 10\,000 trajectoires. Quelle est l'erreur type de ce quantile, en écarts types et en part de la VaR~? Combien de trajectoires faudrait-il pour ramener cette erreur type à 0.5~\%~?
\end{question}

\begin{question}[$\star\star$ \iqroles{Risque, QR} \iqfirm{toutes} \iqlevel{junior, confirmé}]\label{iq:iv2:risk-validation-and-regulation:17}
Pour deux positions dont les volatilités en dollars sont de 0.5 et 0.45 et dont la corrélation est de 0.8, décomposez la VaR à un jour à 99~\% en contributions dont la somme est égale au total. Pourquoi se somment-elles~? Comment le vérifiez-vous par un bump~?
\end{question}

\begin{question}[$\star\star$ \iqroles{Risque, Trader, Analyste} \iqfirm{banque, hedge fund} \iqlevel{junior, confirmé}]\label{iq:iv2:risk-validation-and-regulation:18}
Un portefeuille est long de 50 millions de dollars d'actions de bêta 1, présente un DV01 de 40\,000 dollars en duration longue, et un CS01 de 30\,000 dollars (il perd ce montant lorsque les spreads s'élargissent d'un point de base). Appliquez aux actions $-20~\%$, aux taux $-100$ points de base et aux spreads $+150$. Quel est le P\&L~? De combien les actions doivent-elles baisser davantage pour une perte de 15 millions de dollars~?
\end{question}

\begin{question}[$\star\star$ \iqroles{Risque, QR, Analyste} \iqfirm{hedge fund, banque} \iqlevel{junior, confirmé}]\label{iq:iv2:risk-validation-and-regulation:19}
Trois actions ont pour pondérations 0.5, 0.3 et 0.2, et pour bêtas sur deux facteurs (1.2~; 0.3), (0.8~; $-0.2$) et (1.0~; 0.5). La covariance des facteurs a pour variances 0.04 et 0.01 et pour covariance 0.006. Les variances idiosyncratiques valent 0.03, 0.02 et 0.04. Quelles sont la volatilité du portefeuille et la part de sa variance provenant des facteurs~?
\end{question}

\begin{question}[$\star\star$ \iqroles{Risque, Trader, QT} \iqfirm{hedge fund} \iqlevel{junior, confirmé}]\label{iq:iv2:risk-validation-and-regulation:20}
Vous êtes long de 10 millions de dollars d'une action et en position courte de 10 millions de dollars d'un pair, chacune ayant une volatilité quotidienne de 2~\%. Quelle est la VaR à un jour à 99~\% pour une corrélation de 0.9~? Quelle est-elle si la corrélation tombe à 0.5~? Qu'en dit l'usage d'une corrélation récente~?
\end{question}

\begin{question}[$\star\star$ \iqroles{Risque, QR} \iqfirm{toutes} \iqlevel{junior, confirmé}]\label{iq:iv2:risk-validation-and-regulation:21}
La distribution A est une loi normale centrée réduite. La distribution B est une loi normale centrée réduite avec une probabilité de 99.5~\%, et une perte de 50 avec une probabilité de 0.5~\%. Comparez la VaR à 99~\% et l'expected shortfall à 99~\%. Pourquoi les régulateurs se sont-ils tournés vers l'expected shortfall~? Que perd-on à ce changement~?
\end{question}

\begin{question}[$\star\star$ \iqroles{Risque, Trader, QT} \iqfirm{hedge fund, banque} \iqlevel{junior, confirmé}]\label{iq:iv2:risk-validation-and-regulation:22}
Vous détenez 10 millions de dollars d'une action aux rendements normaux, de volatilité quotidienne de 2~\% et de demi-fourchette de 20 points de base, si bien qu'une vente coûte 20 points de base du montant vendu. Vous devez vendre en cinq parts égales, une à chaque clôture quotidienne, de sorte que la position entière est exposée dès le premier jour. Quelle est la VaR à 99~\% de la vente, la fourchette étant ajoutée comme coût certain~? Quelle est-elle si vous pouvez tout vendre à la première clôture~?
\end{question}

\begin{question}[$\star\star$ \iqroles{Risque, QR, QD} \iqfirm{banque} \iqlevel{junior, confirmé}]\label{iq:iv2:risk-validation-and-regulation:23}
Un desk vous demande de valider son nouveau modèle de pricing avant sa mise en production. Que faites-vous~? Donnez les étapes dans l'ordre, ainsi que ce qui vous conduirait à le refuser.
\end{question}

\begin{question}[$\star\star$ \iqroles{Risque, QR} \iqfirm{banque} \iqlevel{junior, confirmé}]\label{iq:iv2:risk-validation-and-regulation:24}
Une VaR à 99~\% est soumise à un feu tricolore sur 250 jours. Si le modèle est juste, quelle est la probabilité qu'elle reste au vert~? Si le taux réel de dépassement vaut 2~\% ou 3~\%, quelle est la probabilité qu'elle reste au vert, et qu'elle passe au rouge~?
\end{question}

\begin{question}[$\star\star$ \iqroles{Risque, Analyste, QR} \iqfirm{banque} \iqlevel{junior, confirmé}]\label{iq:iv2:risk-validation-and-regulation:25}
Le capital vaut $\max(\mathrm{VaR}_{t-1}, m\,\overline{\mathrm{VaR}}) + \max(\mathrm{sVaR}_{t-1}, m\,\overline{\mathrm{sVaR}})$. La VaR valait 12 la veille et sa moyenne sur 60 jours est de 10. La VaR stressée valait 18 et sa moyenne est de 20, et $m = 3.65$. Quel est le capital~? Jusqu'à quel niveau faudrait-il que la VaR de la veille monte pour compter~? Qu'apporte la VaR stressée~?
\end{question}

\begin{question}[$\star\star\star$ \iqroles{Risque, QR, Trader} \iqfirm{toutes} \iqlevel{confirmé}]\label{iq:iv2:risk-validation-and-regulation:26}
Un portefeuille couvert en delta est vendeur d'options de gamma $-500$ par dollar au carré, et le mouvement quotidien de l'action suit une loi normale d'écart type de 2 dollars. Négligez le thêta. Que dit la VaR delta-normale~? Quelle est la VaR exacte à 99~\%~? Que donne une approximation normale fondée sur la moyenne et la variance~?
\end{question}

\begin{question}[$\star\star\star$ \iqroles{Risque, QR} \iqfirm{toutes} \iqlevel{confirmé}]\label{iq:iv2:risk-validation-and-regulation:27}
Vos rendements présentent une asymétrie de $-0.5$ et un kurtosis en excès de 3. Donnez le quantile à 1~\% de Cornish-Fisher. Utilisez maintenant la formule pour une loi de Student-$t$ à 6 degrés de liberté, normalisée à variance unitaire, dont le quantile à 1~\% du $t_6$ non normalisé vaut $-3.143$. Qu'apprenez-vous~?
\end{question}

\begin{question}[$\star\star\star$ \iqroles{Risque, QR} \iqfirm{banque} \iqlevel{confirmé}]\label{iq:iv2:risk-validation-and-regulation:28}
Un desk annonce un expected shortfall à 97.5~\% sous un modèle normal. Sur $T = 250$ jours, une statistique de test vaut $Z = 1 + \sum_t X_t I_t/(T \times 0.025 \times \mathrm{ES})$, où $X_t$ est le P\&L et $I_t$ marque les jours au-delà de la VaR. Quelle est son espérance si le modèle est juste~? Et si la volatilité réelle vaut 1.2 fois celle du modèle~? Où placeriez-vous le seuil de 5~\%, et quelle puissance a-t-il~?
\end{question}

\begin{question}[$\star\star\star$ \iqroles{Risque, QR} \iqfirm{toutes} \iqlevel{junior, confirmé}]\label{iq:iv2:risk-validation-and-regulation:29}
Une banque utilise une simulation historique sur 250 jours pour une VaR à 99~\%. La volatilité double du jour au lendemain et reste doublée pendant 20 jours. Quel est le taux de dépassements sur ces jours, et quel est le nombre attendu~? Comment se comporte un modèle à pondération exponentielle de paramètre $\lambda = 0.94$ le dixième jour~?
\end{question}

\begin{question}[$\star\star\star$ \iqroles{Risque, QR} \iqfirm{toutes} \iqlevel{confirmé}]\label{iq:iv2:risk-validation-and-regulation:30}
Les pertes au-delà d'un seuil de 2 suivent une loi de Pareto généralisée de paramètre de forme 0.2 et d'échelle 1, et 5~\% des jours dépassent le seuil. Quelles sont la VaR à 99~\% et l'expected shortfall à 99~\%~? Que se passe-t-il lorsque la forme atteint 1~?
\end{question}

\begin{question}[$\star\star\star$ \iqroles{Risque, QR} \iqfirm{toutes} \iqlevel{confirmé}]\label{iq:iv2:risk-validation-and-regulation:31}
Deux actifs ont une corrélation de 0.5. Quelle est la probabilité que les deux passent sous leur propre quantile à 1~\%, sous une copule gaussienne, puis sous une copule de Student-$t$ à 4 degrés de liberté~? Quelle est la limite du coefficient de dépendance de queue dans chaque cas~?
\end{question}

\begin{question}[$\star\star\star$ \iqroles{Risque, QR} \iqfirm{banque} \iqlevel{confirmé}]\label{iq:iv2:risk-validation-and-regulation:32}
Selon les règles relatives au portefeuille de négociation (trading book), une expected shortfall à 97.5~\% est calculée sur un horizon de base de 10 jours, puis étendue en cascade par horizons de liquidité de 10, 20, 40, 60 et 120 jours. L'expected shortfall à 10 jours sur l'ensemble des facteurs vaut 10, et sur les facteurs dont l'horizon est d'au moins 20, 40, 60 et 120 jours, elle vaut respectivement 8, 5, 3 et 2. Quelle est la mesure de capital~? Pourquoi n'est-elle pas $10\sqrt{12}$~?
\end{question}

\begin{question}[$\star\star\star$ \iqroles{Risque, QR, Dev} \iqfirm{toutes} \iqlevel{junior, confirmé}]\label{iq:iv2:risk-validation-and-regulation:33}
Un tableur de VaR calcule une variation de taux comme $(x_t - x_{t-1})/(x_t + x_{t-1})$ alors qu'il devait diviser par la moyenne des deux valeurs. Que deviennent la volatilité et la VaR~? Combien de dépassements par an attendriez-vous dans un backtest, pour des rendements normaux~? Comment un validateur repérerait-il cette erreur avant tout backtest~?
\end{question}

\begin{question}[$\star\star\star$ \iqroles{Risque, QR} \iqfirm{banque} \iqlevel{confirmé}]\label{iq:iv2:risk-validation-and-regulation:34}
L'exposition en cas de défaut (EAD) vaut $1.4\,(\mathrm{RC} + \mathrm{mult} \times \mathrm{AddOn})$, avec $\mathrm{mult} = \min\bigl(1,\ 0.05 + 0.95\exp((V - C)/(2 \times 0.95 \times \mathrm{AddOn}))\bigr)$ et $\mathrm{RC} = \max(V - C, 0)$. Prenez un add-on de 10. Quelle vaut-elle pour $V - C = 5$~? Pour $-20$~? Quel est le plancher~?
\end{question}

\begin{question}[$\star\star\star$ \iqroles{Risque, QR} \iqfirm{banque, hedge fund} \iqlevel{confirmé}]\label{iq:iv2:risk-validation-and-regulation:35}
Le P\&L d'un portefeuille vaut $2x_1 - x_2$ en millions, où les facteurs ont une espérance nulle, des variances de 4 et 1, et une covariance de 1. Quel est le scénario le plus probable qui perd exactement 10~? À combien d'écarts types se trouve-t-il~? Pourquoi le second facteur évolue-t-il comme il le fait~?
\end{question}

\begin{question}[$\star\star\star$ \iqroles{Risque, QR} \iqfirm{toutes} \iqlevel{junior, confirmé}]\label{iq:iv2:risk-validation-and-regulation:36}
Un risk manager rédige une matrice de corrélations stressée pour trois actifs~: corrélations de 0.9 pour les couples (1, 2) et (2, 3), et $-0.5$ pour le couple (1, 3). Qu'est-ce qui ne va pas~? Dans quel intervalle peut varier la corrélation (1, 3)~? Que peut-il mal tourner si vous utilisez malgré tout la matrice, et comment la corrigez-vous~?
\end{question}

\begin{question}[$\star\star\star$ \iqroles{Risque, QR} \iqfirm{banque} \iqlevel{confirmé}]\label{iq:iv2:risk-validation-and-regulation:37}
Vous effectuez un backtest d'une VaR à 99~\% sur un horizon de dix jours, avec des fenêtres de dix jours qui se chevauchent, sur 250 jours. Pour des rendements quotidiens normaux indépendants, quels sont l'espérance et la variance du nombre de dépassements, comparés à une loi binomiale supposant des jours indépendants~? Quel effet cela a-t-il sur le test usuel, et que faites-vous pour y remédier~?
\end{question}

\begin{question}[$\star\star\star\star$ \iqroles{Risque, QR} \iqfirm{banque} \iqlevel{confirmé}]\label{iq:iv2:risk-validation-and-regulation:38}
Vous validez la VaR à 99~\% à un jour du desk. Sur 500 jours, elle a enregistré 9 dépassements. Les transitions sont~: 485 non-non, 5 non-oui, 5 oui-non, 4 oui-oui. Lancez les deux tests. Puis passez en revue ce que vous regarderiez ensuite, et ce que vous diriez au desk et au régulateur.
\end{question}

\begin{question}[$\star\star\star\star$ \iqroles{Risque, Trader, QR} \iqfirm{banque, hedge fund} \iqlevel{confirmé}]\label{iq:iv2:risk-validation-and-regulation:39}
Un book d'options est couvert en delta, avec un gamma de $-20\,000$ actions par dollar, un vega de $-\$40\,000$ par point de volatilité, et un theta de $+\$22\,500$ par jour. Le mouvement quotidien de l'action suit une loi normale d'écart type de 1.50~\$. La volatilité implicite évolue selon une loi normale de 1 point par jour, avec une corrélation de $-0.7$ avec l'action. Estimez la VaR et l'expected shortfall à 99~\%~: delta-normale, gamma seul, gamma et vega ensemble. Que se passe-t-il si les mouvements suivent une loi de Student-$t$ à 4 degrés de liberté~? Qu'en est-il d'un choc de $-\$10$ et de $+5$ points~? Quelle limite fixez-vous~?
\end{question}

\begin{question}[$\star\star\star\star$ \iqroles{Risque, QR, QD} \iqfirm{banque} \iqlevel{confirmé}]\label{iq:iv2:risk-validation-and-regulation:40}
L'arbre binomial d'un collègue pour un call européen donne les résultats suivants, avec $S = K = 100$, $r = 5~\%$, $\sigma = 20~\%$ et $T = 1$. La formule fermée vaut 10.4506. Pour $N = 50$, 10.4107. Pour 51, 10.4850. Pour 100, 10.4306. Pour 101, 10.4680. Pour 200, 10.4406. Pour 201, 10.4593. Pour un strike de 180, $N = 50$ donne 0.0249 et $N = 200$ donne 0.0276, et la formule fermée vaut 0.0286. Pour $\sigma = 0.1~\%$ et $N = 500$, il renvoie $2.7 \times 10^{159}$. Que concluez-vous~? Quelle est la correction pour chacun~? Que testez-vous d'autre avant de signer~?
\end{question}
